\documentclass[11pt]{amsart}

\usepackage{amsmath,amssymb,amsthm,mathtools}
\usepackage{enumitem}
\usepackage[hidelinks]{hyperref}
\usepackage{microtype}

% Paths relative to this manuscript directory
% PeAIce claim-discipline macros (input from manuscript preambles)
% Status tags for theorem environments and inline use

\newcommand{\PeAIceFormal}{\textsc{Formal}}
\newcommand{\PeAIceProposed}{\textsc{Proposed}}
\newcommand{\PeAIceOpen}{\textsc{Open}}
\newcommand{\PeAIceKnown}{\textsc{Known}}
\newcommand{\PeAIceNumerics}{\textsc{Numerics}}
\newcommand{\PeAIceClosedNeg}{\textsc{Closed-Negative}}
\newcommand{\PeAIceClosedPos}{\textsc{Closed-Positive}}
\newcommand{\PeAIceAnalogy}{\textsc{Structural Analogy}}

\newcommand{\statusbox}[1]{%
  \par\noindent\fbox{\parbox{0.95\linewidth}{\small\textbf{Status.} #1}}\par\medskip
}

\newcommand{\firewallbox}[1]{%
  \par\noindent\fbox{\parbox{0.95\linewidth}{\small\textbf{Firewall.} #1}}\par\medskip
}

\newcommand{\PeAIceOpenTargets}{%
  The {R}iemann hypothesis and the {C}oleman conjecture remain \PeAIceOpen.
  No claim in this note constitutes a proof of either statement.
}

\newcommand{\ZetaFirewall}{%
  The value $\zeta(0)=-\tfrac12$ is \emph{not} asserted as a term in the
  {K}akeya proof literature. Any regularization at $s=0$ is second-stage and
  conditional on a residual operator or counting object constructed after
  residual control.
}

\newcommand{\KNSFirewall}{%
  Typed incidence objects do not locate zeros of $\zeta$.
  Overlap measures are not placement measures on the critical line.
}


\theoremstyle{plain}
\newtheorem{theorem}{Theorem}[section]
\newtheorem{lemma}[theorem]{Lemma}
\newtheorem{proposition}[theorem]{Proposition}
\newtheorem{corollary}[theorem]{Corollary}

\theoremstyle{definition}
\newtheorem{definition}[theorem]{Definition}
\newtheorem{remark}[theorem]{Remark}
\newtheorem{conjecture}[theorem]{Conjecture}

\newcommand{\dd}{\mathrm{d}}
\newcommand{\Gzero}{G_0}
\newcommand{\Cdr}{C_{\delta,\rho}}
\newcommand{\Tdelta}{\mathcal{T}_\delta}
\newcommand{\edelta}{e_\delta}
\newcommand{\mdelta}{m_\delta}

\title[Grain Zero residual]{Grain Zero: residual overlap after structured
{K}akeya factorization in $\mathbb{R}^3$}
\author{Manuel Coleman}
\address{Love Labs LCA / PeAIce Research Program}
\email{coleman.manuel23@gmail.com}
\date{\today}

\begin{document}

\begin{abstract}
We isolate a residual overlap measure---\emph{Grain Zero}---after a finite
family of $\delta$-tubes in $\mathbb{R}^3$ has been factored into the
structured carriers appearing in modern proofs of the Kakeya set conjecture
in three dimensions (grains, slabs, prisms, planks, sticky packets, convex
carriers, and scale envelopes). The raw excess multiplicity is not itself
Grain Zero; Grain Zero is the excess supported off the extracted carrier.
The primary theorem target is residual vanishing: after the Wang--Zahl /
Guth--Wang--Zahl reduction stages, the normalized residual mass should tend
to zero as $\delta\to 0$. Operator encodings and any zeta-regularization near
$s=0$ are second-stage and conditional.
\PeAIceOpenTargets\
\ZetaFirewall
\end{abstract}

\subjclass[2020]{42B25, 28A75, 42B20}
\keywords{Kakeya set conjecture, residual measures, incidence geometry,
multiplicity, Wang--Zahl reduction}

\maketitle

\section{Introduction}

\statusbox{Program note ARX-001 under PeAIce claim discipline.
Definition of $\Gzero$ is \PeAIceFormal.
Residual vanishing is \PeAIceProposed/\PeAIceOpen.
Operator and zeta stages are explicitly second-stage.}

A Kakeya (Besicovitch) set in $\mathbb{R}^n$ contains a unit line segment in
every direction. In $\mathbb{R}^3$, every such set has Minkowski and Hausdorff
dimension three \cite{WangZahl2025Kakeya,GuthWangZahl2026Streamlined}. The
proof architecture organizes directional overlap into structured carriers.
This note does not reprove the Kakeya theorem. It asks a residual question:
after legal structure is removed, what measure remains, and can it be shown
to vanish under the same estimates that close the volume problem?

We call the residual object \emph{Grain Zero}. The name is program-internal;
the mathematics is a residual measure after carrier subtraction.

\firewallbox{\ZetaFirewall}

\section{Tube setup}

Let $\Tdelta$ be a finite family of $\delta$-tubes in $\mathbb{R}^3$ with
shading $Y(T)\subset T$ for each $T\in\Tdelta$. Write
\begin{equation}
  U(\Tdelta,Y) \;:=\; \bigcup_{T\in\Tdelta} Y(T),
\end{equation}
\begin{equation}
  \mdelta(x) \;:=\; \#\{\,T\in\Tdelta : x\in Y(T)\,\},
\end{equation}
\begin{equation}
  \edelta(x) \;:=\; \bigl(\mdelta(x)-1\bigr)_+.
\end{equation}
The raw excess measure $\edelta(x)\,\dd x$ still includes overlap that may be
fully accounted for by structured carriers. Grain Zero begins only after that
structure is extracted.

\section{Structured carriers}

\begin{definition}[Structured carrier, informal]
\label{def:carrier}
At a scale passage $\delta\le \rho\le 1$, let $\Cdr$ denote the (measurable)
structured carrier produced by a fixed stage of a Wang--Zahl /
Guth--Wang--Zahl style reduction. Depending on the stage, $\Cdr$ may include
grains, slabs, prisms, planks, convex carriers, sticky packets, Katz--Tao
convex-density packets, Frostman slab packets, high-multiplicity cells, and
scale-induction envelopes.
\end{definition}

\begin{remark}
A rigorous implementation must pin $\Cdr$ to explicit source addresses in
\cite{WangZahl2025Kakeya,GuthWangZahl2026Streamlined} for the chosen stage.
The present scaffold records the program definition; carrier pinning is an
owed edit before submission.
\end{remark}

\section{Grain Zero}

\begin{definition}[Grain Zero residual measure]
\label{def:g0}
After $\Cdr$ has been extracted, define
\begin{equation}
  \dd \Gzero_{\delta,\rho}(x)
  \;:=\;
  \edelta(x)\, \mathbf{1}_{\mathbb{R}^3\setminus \Cdr}(x)\, \dd x.
\end{equation}
Equivalently, for measurable $A\subset\mathbb{R}^3$,
\begin{equation}
  \Gzero_{\delta,\rho}(A)
  \;:=\;
  \int_A \bigl(\mdelta(x)-1\bigr)_+
  \,\mathbf{1}_{\mathbb{R}^3\setminus \Cdr}(x)\, \dd x.
\end{equation}
\end{definition}

\begin{remark}
$\Gzero$ is not a grain, slab, prism, or sticky packet. It is the leftover
excess after those carriers absorb the overlap they are designed to control.
\end{remark}

\section{Primary theorem target}

\begin{conjecture}[Residual vanishing]
\label{conj:vanish}
Under the hypotheses and refinements of the Wang--Zahl /
Guth--Wang--Zahl reduction, after broad/narrow decomposition, convex-density
factoring, grain/slab/prism extraction, sticky or non-sticky reduction, and
final scale refinement,
\begin{equation}
  \frac{\Gzero_{\delta,\rho}(\mathbb{R}^3)}
       {\sum_{T\in\Tdelta} |Y(T)|}
  \;\longrightarrow\; 0
  \qquad\text{as }\delta\to 0.
\end{equation}
\end{conjecture}

\statusbox{Conjecture~\ref{conj:vanish} is \PeAIceProposed/\PeAIceOpen.
A weak form with $\delta^{o(1)}$ capture may be assemblable from existing
error channels; a strong form is the program's primary geometric burden.}

\section{Lemma ladder (program outline)}

We record the intended reduction chain. Items not proved here are marked.

\begin{enumerate}[label=\textbf{L\arabic*.},leftmargin=*]
  \item \textbf{Structured factorization} (\PeAIceProposed).
  Decompose $\edelta\,\dd x = \dd S_{\delta,\rho} + \dd \Gzero_{\delta,\rho}$
  with $\dd S$ supported on carriers and $\dd \Gzero$ off carriers.
  \item \textbf{Carrier exhaustion} (\PeAIceProposed).
  Show $\dd S$ accounts for overlap controlled by existing incidence,
  multiplicity, convex-density, or slab estimates at the chosen stage.
  \item \textbf{Residual domination} (\PeAIceProposed).
  Bound
  \[
    \Gzero_{\delta,\rho}(\mathbb{R}^3)
    \le
    E_{\mathrm{bn}} + E_{\mathrm{convex}} + E_{\mathrm{sticky}} + E_{\mathrm{scale}}.
  \]
  \item \textbf{Residual vanishing} (\PeAIceOpen).
  Deduce Conjecture~\ref{conj:vanish}.
  \item \textbf{Operator encoding} (second stage, \PeAIceOpen).
  Only after residual control, construct $K_{\Gzero}$, a kernel, a counting
  function, or a scale defect measure from residual data.
  \item \textbf{Zeta-regularization question} (conditional, \PeAIceOpen).
  Only after an operator exists, ask whether a zeta/Mellin object continues
  near $s=0$. \ZetaFirewall
\end{enumerate}

\section{Obstacles}

\begin{enumerate}[leftmargin=*]
  \item $\Cdr$ is proof-stage dependent; several residual lanes
  $\Gzero^{\mathrm{grain}},\Gzero^{\mathrm{slab}},\ldots$ may be required.
  \item Existing Kakeya proofs close by volume/multiplicity control without
  naming a residual measure; residual closure must be \emph{translated} from
  those estimates.
  \item A vanishing residual may be too thin to support a nontrivial spectral
  object; the correct second-stage object may be a renormalized defect along
  scales rather than $\Gzero$ itself.
\end{enumerate}

\section{Related work and non-claims}

The geometric background is the Kakeya theorem in $\mathbb{R}^3$
\cite{WangZahl2025Kakeya,GuthWangZahl2026Streamlined} and the broader
restriction/Kakeya complex \cite{BennettCarberyTao2006,Guth2025LargeValues}.
We do not assert a peer-reviewed literature bridge from Kakeya residual
measures to zeros of the Riemann zeta function; any such bridge is
program-original and presently \PeAIceOpen.

\PeAIceOpenTargets

\section*{Acknowledgments}
Research engineering assistance was used in preparing program materials.
Mathematical claims and any future arXiv submission decision remain the
author's responsibility. Evaluator non-sovereignty ($h<1$) is a governance
constraint on the program, not a theorem of analytic number theory.

\bibliographystyle{amsplain}
\bibliography{../../bibliography/peaice-arxiv}

\end{document}
